\slajdid{09_1-s056}
\begin{frame}[label=09_1-s056]{Problem koji nismo razmatrali}
\gusttekst\setlength{\abovedisplayskip}{3pt}\setlength{\belowdisplayskip}{3pt}
\begin{columns}[c,onlytextwidth]\column{.24\textwidth}\centering\begin{tikzpicture}[x=.8cm,y=.65cm,font=\sitneoznake,>=Latex]\begin{scope}[shift={(0,1.4)}]\draw(-.9,-.2)--(-.15,-.2);\draw(-.15,-.34)--(-.15,.34);\foreach\a/\b in{-.4/-.22,-.1/.1,.22/.4}{\draw(0,\a)--(0,\b);}\draw(0,.3)--(.4,.3)--(.4,.7);\draw(0,-.3)--(.4,-.3)--(.4,-.7);\draw[-{Triangle[length=2mm,width=1.4mm]}](.75,0)--(.03,0);\node[left]at(-.9,-.2){$V_C$};\draw[fill=white](-.9,-.2)circle(.05);\node[right]at(.6,.35){$T_L$};\draw(.75,0)--(.75,-.3)--(.4,-.3);\end{scope}\begin{scope}[shift={(0,0)}]\draw(-.9,-.2)--(-.15,-.2);\draw(-.15,-.34)--(-.15,.34);\foreach\a/\b in{-.4/-.22,-.1/.1,.22/.4}{\draw(0,\a)--(0,\b);}\draw(0,.3)--(.4,.3)--(.4,.7);\draw(0,-.3)--(.4,-.3)--(.4,-.7);\draw[-{Triangle[length=2mm,width=1.4mm]}](.75,0)--(.03,0);\node[left]at(-.9,-.2){$V_I$};\draw[fill=white](-.9,-.2)circle(.05);\node[right]at(.6,.35){$T_D$};\draw(.75,0)--(.75,-.3)--(.4,-.3);\end{scope}\draw(.2,2.1)--(.6,2.1)node[above left]{$V_{DD}$};\draw(.4,-.7)--(.4,-.9)node[ground]{};\draw(.4,.7)--(1.4,.7)node[ocirc]{}node[right]{$V_O$};\draw[DEplava](-1.15,-1)--(1.6,2.5);\draw[DEplava](-1.15,2.5)--(1.6,-1);\end{tikzpicture}
\column{.06\textwidth}\centering\tikz{\draw[->,DEplava](0,0)--(.75,0);}\column{.24\textwidth}\centering\begin{tikzpicture}[x=.8cm,y=.65cm,font=\sitneoznake,>=Latex]\begin{scope}[shift={(0,1.4)}]\draw(-.9,-.2)--(-.15,-.2);\draw(-.15,-.34)--(-.15,.34);\foreach\a/\b in{-.4/-.22,-.1/.1,.22/.4}{\draw(0,\a)--(0,\b);}\draw(0,.3)--(.4,.3)--(.4,.7);\draw(0,-.3)--(.4,-.3)--(.4,-.7);\draw[-{Triangle[length=2mm,width=1.4mm]}](.75,0)--(.03,0);\node[left]at(-.9,-.2){$V_C$};\draw[fill=white](-.9,-.2)circle(.05);\node[right]at(.6,.35){$T_L$};\draw(.75,0)--(.95,0);\end{scope}\begin{scope}[shift={(0,0)}]\draw(-.9,-.2)--(-.15,-.2);\draw(-.15,-.34)--(-.15,.34);\foreach\a/\b in{-.4/-.22,-.1/.1,.22/.4}{\draw(0,\a)--(0,\b);}\draw(0,.3)--(.4,.3)--(.4,.7);\draw(0,-.3)--(.4,-.3)--(.4,-.7);\draw[-{Triangle[length=2mm,width=1.4mm]}](.75,0)--(.03,0);\node[left]at(-.9,-.2){$V_I$};\draw[fill=white](-.9,-.2)circle(.05);\node[right]at(1.02,.25){$T_D$};\draw(.75,0)--(.95,0);\end{scope}\draw(.2,2.1)--(.6,2.1)node[above left]{$V_{DD}$};\draw(.4,-.7)--(.4,-.9)node[ground]{};\draw(.95,1.4)--(.95,-.9)node[ground]{};\draw(.4,.7)--(.83,.7);\draw[preaction={draw=white,line width=2.5pt}](.83,.7)arc[start angle=180,end angle=0,radius=.12];\draw(1.07,.7)--(1.4,.7)node[ocirc]{}node[right]{$V_O$};\end{tikzpicture}
\column{.44\textwidth}
\[V_{Tn}=V_{Tn0}+\gamma\left(\sqrt{|2\phi_F+V_{SB}|}-\sqrt{|2\phi_F|}\right)\]
\[V_{Tn,L}=V_{Tn0}+\gamma\left(\sqrt{|2\phi_F+V_{SB,L}|}-\sqrt{|2\phi_F|}\right)\]
\end{columns}
Pri čemu je
\[V_{SB,L}=V_{S,L}-V_{B,L}=V_{OH}-0=V_{OH}\]
Znači
\[V_{OH}=V_{DD}-V_{Tn0}-\gamma\left(\sqrt{|2\phi_F+V_{OH}|}-\sqrt{|2\phi_F|}\right)\]
i to treba, i može da se reši, po $V_{OH}$. Bilo računanjem kvadratne jednačine, bilo iterativno. Ono što je uočljivo jeste da će napon $V_{OH}$ značajno pasti. Slično važi i za tačku $V_{OL}$. Tačka $V_{IL}$ ne zavisi od $V_{Tn,L}$ ostaje ista. Na žalost ovo ne važi za tačku i $V_{IH}$ pošto postoji zavisnost napona praga od napona sors – osnovna, odnosno od izlaznog napona
\end{frame}
